\documentclass[11pt]{article}

\usepackage[T1]{fontenc}
\usepackage[utf8]{inputenc}
\usepackage{amsmath,amssymb,amsthm}
\usepackage{array,booktabs}
\usepackage[margin=1in]{geometry}
\usepackage[hidelinks]{hyperref}

\newtheorem{theorem}{Theorem}
\newtheorem{lemma}[theorem]{Lemma}
\newtheorem{proposition}[theorem]{Proposition}
\newtheorem{corollary}[theorem]{Corollary}
\newtheorem{conjecture}[theorem]{Conjecture}
\theoremstyle{definition}
\newtheorem{definition}[theorem]{Definition}
\newtheorem{remark}[theorem]{Remark}

\newcommand{\DTIME}{\mathsf{DTIME}}
\newcommand{\DTISP}{\mathsf{DTISP}}
\newcommand{\NTIME}{\mathsf{NTIME}}
\newcommand{\SAT}{\mathsf{SAT}}
\newcommand{\leanus}{\textunderscore\allowbreak}
\newcommand{\lean}[1]{\ifmmode\text{\ttfamily #1}\else{\ttfamily\let\_\leanus #1}\fi}
\newcommand{\pin}{\mathrm{pin}}
\newcommand{\cstair}{c^{*}_{\mathrm{stair}}}

\title{Williams's alternation-trading calculus for SAT on $d$-dimensional one-tape machines:\\ the optimality conjecture fails, and the threshold lies in a short interval}
\author{PastaPastaPasta}
\date{November 2026}

\emergencystretch=2em

\begin{document}
\maketitle

\begin{abstract}
Williams \cite[\S7]{W13} sets up a calculus of alternation-trading proofs for lower bounds against machines with random-access input, $n^{o(1)}$ bits of random-access storage and one $d$-dimensional sequential work tape. His Theorem~7.5 proves that SAT is not solvable in time $n^c$ on such machines for every $c<r_d$, where $r_d$ is a root of a quintic ($r_1\approx1.30094$, $r_2\approx1.18879$, $r_3\approx1.13728$). He conjectures that $r_d$ is optimal for alternation-trading proofs, and proves (in a sketch, Theorem~7.9) that no proof in the calculus reaches $c\ge1+1/(d+1)$.

We study a Lean~4 formalisation of this calculus. For $d=1,2,3$ it has normal-form proofs at $c=1.30921174$, $1.1915968$ and $1.1386523$, all above $r_d$, so the conjecture fails for the formalised calculus. These proofs also satisfy a stricter reading of the speedup lemma (kernel time at least the input length, and space at most time, on every line). Conversely, for every $c\ge1.3094$, $1.192$ and $1.1392$ respectively there is no normal-form proof, strict or not, even in a relaxed calculus that contains every move of Williams's calculus as we read it (Section~\ref{sec:coverage} discusses one sound move outside it). So the threshold of the formalised calculus lies in $[1.30921174,\,1.3094]$, $[1.1915968,\,1.192]$ and $[1.1386523,\,1.1392]$. The constructions are shared-subproof certificates checked by a verified checker; the impossibility results are multi-level potential functions whose Lemma~7.1 condition is checked for all words at once by a finite abstraction. The exact threshold is open. A recursion read off from the optimal proofs has threshold $\cstair(1)=1.3092117446776\ldots$; that this is the threshold of the calculus is a conjecture, supported by numerics only.

\emph{What is kernel-checked:} the existence of the proofs at the three lower values, the absence of proofs at and above the three upper values, and the simpler ceiling $c^2\ge1+1/d$, all for the formal calculus. \emph{What is not:} that the formal rules are Williams's (a modelling claim, compared rule by rule in Table~\ref{tab:rules} with the reading choices listed in Section~\ref{sec:reading}), and the soundness of Williams's lemmas. The complexity consequence, $\SAT\notin\DTIME_d[n^c]$ for every $c<1.30921174$, $1.1915968$, $1.1386523$, holds exactly as far as Williams's Lemmas~7.1, 7.2, 7.3, his Lemma~3.3 used inside a prefix, and Theorem~3.1 do. These are the results his Theorem~7.5 rests on, and they are not formalised here. All statements concern one proof calculus.
\end{abstract}

\section{Introduction}

\paragraph{Background.}
Alternation-trading proofs derive a contradiction from a hypothesis such as $H_c$: $\NTIME[n]\subseteq\DTIME_d[n^c]$ by alternately applying \emph{speedup} rules, which add quantifier blocks to a class and reduce its running time, and \emph{slowdown} rules, which use the hypothesis to remove a block at the cost of a factor $c$ in the exponent, until a time hierarchy is contradicted. Williams \cite{W13} formalises such proofs for three settings and studies them by linear programming. Its \S7 treats the ``hybrid'' model: $\DTIME_d[t]$ is the class of languages recognised in time $O(t)$ by a machine with a read-only random-access input tape, $n^{o(1)}$ bits of read-write random-access storage, and one $d$-dimensional read-write tape with sequential access \cite[p.~A:29]{W13}. The speedups are Lemma~7.1 (from $\DTIME_d$ to $\DTISP$, by guessing crossing sequences) and the $\DTISP$ speedup Lemma~3.3; the slowdown is Lemma~7.2.

Theorem~7.5 of \cite{W13} gives $\SAT\notin\DTIME_d[n^c]$ for $c<r_d$, where $r_d\ge1$ is a root of
\[p_d(x)=(2d+1)(d+1)^2x^5-2(d+1)(2d+1)x^4-d^2x^3+2(d+1)x^2-((2d+1)(d+1)^2+1)x+d(d+1).\]
Numerically $r_1=1.3009428844$, $r_2=1.1887894930$, $r_3=1.1372807069$ (script \texttt{staircase.py roots}). After the proof, \cite[p.~A:35]{W13} writes: ``We again conjecture that the above lower bound is optimal for alternation-trading proofs'', and Theorem~7.9 states that no alternation-trading proof gives $c\ge1+1/(d+1)$, with a proof sketch (pp.~A:35--A:37). The long version \cite{Wfull} treats $d=1$ and says ``we feel safe in conjecturing that our above theorem is optimal for alternation-trading proofs''.

Page and result numbers refer to the TOCT author copy of \cite{W13}. Our own results are numbered $1,2,3,\dots$ Every theorem below is a statement about the formal calculus, proved by the Lean theorems named with it; names refer to the accompanying Lean library (namespace \lean{PastaMath.OneTapeThreshold}). The conditional complexity consequence, Corollary~\ref{cor:lb}, is a paper argument.

\paragraph{Results.}
\begin{itemize}
\item \textbf{Proofs above $r_d$} (Theorem~\ref{thm:lower}; \lean{conjecture\_fails}, \lean{HStair.stairD\{1,2,3\}\_snf}). The formal calculus \lean{HY.Step} has strict normal-form proofs at $c=1.30921174$ ($d=1$), $1.1915968$ ($d=2$) and $1.1386523$ ($d=3$). Each has kernel time $\ge$ input superscript and space $\le$ time on every line. Every real root of $p_d$ is below the respective $c$. A 39-line proof at $c=1.30095>r_1$ is checked separately (Theorem~\ref{thm:39}).
\item \textbf{Ceilings} (Theorem~\ref{thm:ceiling}; \lean{HLM.ceiling\_d1}, \lean{HLM.ceiling\_d2}, \lean{HLM.ceiling\_d3}). For every real $c\ge1.3094$, $1.192$, $1.1392$ respectively there is no normal-form proof, strict or non-strict, in \lean{HY.Step} or in the relaxed calculus \lean{HL.RStep} of Section~\ref{sec:relaxed}.
\item \textbf{A closed-form ceiling} (Theorem~\ref{thm:sqrt}; \lean{HL.no\_nf}). For every real $d>0$ and $c\ge1$ with $c^2\ge1+1/d$ there is no strict normal-form proof. Since $\sqrt{1+1/d}<1+1/(d+1)$ for $d\ge1$, this improves the bound of \cite[Theorem~7.9]{W13} for the formal calculus for every $d$.
\item \textbf{The threshold} (Corollary~\ref{cor:interval}; \lean{threshold\_bounds}). The supremum $c^*_d$ of the $c$ with a strict normal-form proof satisfies $c^*_1\in[1.30921174,1.3094]$, $c^*_2\in[1.1915968,1.192]$, $c^*_3\in[1.1386523,1.1392]$. The exact values are open (Section~\ref{sec:threshold}).
\end{itemize}

\begin{table}[ht]
\centering
\small
\begin{tabular}{ccccccc}
\toprule
$d$ & $r_d$ & proof at & no proof for $c\ge$ & $\sqrt{1+1/d}$ & $1+\frac1{d+1}$ & $\cstair(d)$ \\
\midrule
1 & 1.3009429 & 1.30921174 & 1.3094 & 1.4142 & 1.5 & 1.30921174468 \\
2 & 1.1887895 & 1.1915968 & 1.192 & 1.2247 & 1.3333 & 1.19159683545 \\
3 & 1.1372807 & 1.1386523 & 1.1392 & 1.1547 & 1.25 & 1.13865239000 \\
\bottomrule
\end{tabular}
\caption{Exponents for the formal calculus. Column~2 is the exponent of \cite[Theorem~7.5]{W13}; columns~3 and~4 are Theorems~\ref{thm:lower} and~\ref{thm:ceiling}, column~5 is Theorem~\ref{thm:sqrt} (all kernel-checked); column~6 is \cite[Theorem~7.9]{W13}; the last column is the conjectured threshold (Section~\ref{sec:threshold}). Columns~2 and~7 come from \texttt{scripts/staircase.py}.}
\label{tab:summary}
\end{table}

\paragraph{Where the improvement comes from.}
The proofs of \cite[Theorem~7.5]{W13} have annotations $[1\,(1\,1\,0)^k\,0\,(1\,1\,0)^\ell\,0]$ (1 a speedup, 0 a slowdown): under an outer $\exists$ block they make one excursion to a $\forall$ level, where simple rounds run (the inner pattern of Lemma~7.7 and Corollary~7.8 of \cite{W13}), return to the $\exists$ level for further simple rounds, and remove the $\exists$ block. The 39-line proof of Section~\ref{sec:39} makes a second excursion under the same $\exists$ block: the first new block of the second excursion merges into the $\exists$ block, so it costs no alternation. Its annotation has three occurrences of $0\,0$; \cite[\S7.2]{W13} reports that a heuristic search for annotations with three occurrences found none better. Iterating the idea gives proofs whose level structure is a staircase (Section~\ref{sec:stair}). Every step of every proof is an application of Lemma~7.1, Lemma~3.3 or Lemma~7.2 of the kinds that occur in the 25-line proof printed in \cite[\S7.2]{W13}: Lemma~7.1 merged into an existing block, Lemma~3.3 inside a prefix, and Lemma~7.2 under a prefix with input superscript above 1. The proofs are much longer: up to about $7.6\cdot10^{21}$ lines when unfolded, a constant independent of $n$.

\paragraph{Scope.}
All results concern one proof calculus. The complexity reading needs the soundness of Williams's lemmas, which we do not formalise (Section~\ref{sec:consequence}); it rests on the same lemmas as \cite[Theorem~7.5]{W13}. The ceilings say nothing about lower bounds proved by other means.

\paragraph{Novelty.}
We are not aware of an improvement of the exponents $r_d$ after \cite{W13}, of a resolution of its \S7 conjecture, or of a limit for this calculus below $1+1/(d+1)$. Our search was not exhaustive.

\paragraph{AI assistance.}
The proofs, the formalisation and the scripts were produced with Anthropic's Claude models under the author's direction. The kernel-checked Lean proofs are the primary evidence for the theorems about the formal calculus; this text explains them and states what they do and do not establish.

\section{The calculus as formalised}\label{sec:calculus}

\subsection{Lines}
A \emph{line} is a class $(Q_1\,n^{g_1})^{o_1}\cdots(Q_m\,n^{g_m})^{o_m}\,K$ with $Q_i\in\{\exists,\forall\}$ and kernel $K=\DTIME_d[n^a]$ or $K=\DTISP[n^a,n^e]$ \cite[\S2, \S7.1]{W13}. Here $g_i$ is the guess exponent of block $i$ (with $0$ standing for a $\log n$ block) and $o_i$ is its output superscript: block $i$ passes $n^{o_i+o(1)}$ bits to the next stage (Williams's $b_{i+1}$). The \emph{input superscript} of the kernel, $b=\pin$, is $o_m$, or $1$ if $m=0$. In Lean (\lean{HY.Line}) a line is a block list, innermost block first, and a kernel \lean{T a} or \lean{S a e}; \lean{HY.pin} is the input superscript and \lean{HY.Alt} says adjacent blocks have different quantifiers. Exponents are real numbers.

\subsection{Rules}\label{sec:rules}
Table~\ref{tab:rules} lists Williams's rules and their formal versions. In each rule $b$ is the input superscript of the line it is applied to, and $c$ the exponent of the hypothesis $H_c$.

\begin{table}[ht]
\centering\small
\begin{tabular}{>{\raggedright\arraybackslash}p{2.5cm}>{\raggedright\arraybackslash}p{6.4cm}>{\raggedright\arraybackslash}p{3.2cm}>{\raggedright\arraybackslash}p{2.2cm}}
\toprule
Williams & statement as formalised (exact outputs) & \lean{HY.Step} & \lean{HL.RStep} \\
\midrule
Lemma 7.1 (p.~A:29) & $Q^b\DTIME_d[n^a]\subseteq{}$ $Q^b(Q'\,n^{a-s})^{\max(a-s,b)}\DTISP[n^a,n^{ds}]$, $0<s\le a$, $b\ge1$; $Q'$ is merged into the innermost block when it has the same type (\S2, p.~A:5), giving guess $\max(g,a-s)$ and output $\max(a-s,o)$ & \lean{convMerge}; \lean{convNew} (new block of the other type, or the first block) & \lean{convM}, \lean{convN} \\
Lemma 3.3 (p.~A:6) inside a prefix, Speedup Rule 1 (p.~A:9) & $(Q\,n^g)^o\DTISP[n^a,n^e]\subseteq{}$ $(Q\,n^{\max(g,x+e)})^{\max(o,x+e)}$ $(\bar Q\log n)^{\max(o,e)}$ $\DTISP[n^{a-x},n^e]$, $a\ge1$, $0\le e\le a-x$, $0\le x\le a-1$, $o\ge1$ & \lean{speedMerge} & \lean{speedM} \\
Speedup Rule 2 (p.~A:9; unnecessary by p.~A:31) & two new blocks $(Q\,n^{x+e})^{\max(b,x+e)}(\bar Q\log n)^{\max(b,e)}$, $Q$ opposite to the innermost block, same side conditions & \lean{speedNew} & \lean{speedN} \\
Lemma 7.2 (p.~A:30) & $Q^{b_k}(Q'\,n^{a_k})^{b_{k+1}}K[n^{a_{k+1}}]\subseteq{}$ $Q^{b_k}\DTIME_d[n^{c\max(b_k,b_{k+1},a_k,a_{k+1})}]$ for $K=\DTIME_d$ or $\DTISP$; $b_k,b_{k+1},a_{k+1}\ge1$, $a_k\ge0$ & \lean{slow} & \lean{slow} \\
Normal form (p.~A:31), Lemma~7.3, Thm.~7.4 & $\DTIME_d[n^a]\subseteq\cdots\subseteq{}$ $\DTIME_d[n^{a'}]$ & \lean{HY.NF} (strict, $1\le a'<a$) & Thm.~\ref{thm:ceiling}(2) \\
\bottomrule
\end{tabular}
\caption{Williams's rules \cite{W13} and the formal rules. \lean{HY.Step} is defined in \texttt{Calculus.lean}, \lean{HL.RStep} in \texttt{Relaxed.lean} (Section~\ref{sec:relaxed}).}
\label{tab:rules}
\end{table}

Precisely, \lean{HY.Step c d} has the constructors
\begin{itemize}
\item \lean{convMerge u rest a s}: from $(u::\mathit{rest},\ \lean{T}\,a)$ with $0<s\le a$, $1\le u.o$, to $(\langle u.q,\max(u.g,a-s),\max(a-s,u.o)\rangle::\mathit{rest},\ \lean{S}\,a\,(ds))$;
\item \lean{convNew bl q a s}: if the innermost block of $\mathit{bl}$ (if any) has quantifier $\ne q$, $0<s\le a$ and $1\le\pin(\mathit{bl})$, from $(\mathit{bl},\lean{T}\,a)$ to $(\langle q,a-s,\max(a-s,\pin\,\mathit{bl})\rangle::\mathit{bl},\ \lean{S}\,a\,(ds))$;
\item \lean{speedMerge u rest a e x}: if $1\le a$, $0\le e\le a-x$, $0\le x\le a-1$ and $1\le u.o$, from $(u::\mathit{rest},\lean{S}\,a\,e)$ to $(\langle\overline{u.q},0,\max(u.o,e)\rangle::\langle u.q,\max(u.g,x+e),\max(u.o,x+e)\rangle::\mathit{rest},\ \lean{S}\,(a-x)\,e)$;
\item \lean{speedNew bl q a e x}: if the innermost block of $\mathit{bl}$ (if any) has quantifier $\ne q$, the same conditions on $a,e,x$ hold and $1\le\pin\,\mathit{bl}$, from $(\mathit{bl},\lean{S}\,a\,e)$ to $(\langle\bar q,0,\max(\pin\,\mathit{bl},e)\rangle::\langle q,x+e,\max(\pin\,\mathit{bl},x+e)\rangle::\mathit{bl},\ \lean{S}\,(a-x)\,e)$;
\item \lean{slow u rest k}: if $1\le\pin\,\mathit{rest}$, $1\le u.o$, $1\le\mathrm{time}(k)$ and $0\le u.g$, from $(u::\mathit{rest},k)$ to $(\mathit{rest},\ \lean{T}\,(c\max(\pin\,\mathit{rest},u.o,u.g,\mathrm{time}(k))))$.
\end{itemize}

\subsection{Normal form}
\begin{definition}[\lean{HY.NF c d M n}]\label{def:nf}
A \emph{strict normal-form proof} at $c$ in dimension $d$ is a list of lines $M_0,\dots,M_n$ with $n\ge1$, each obtained from the previous one by a \lean{HY.Step c d}, such that $M_0=\DTIME_d[n^a]$ and $M_n=\DTIME_d[n^{a'}]$ have no blocks, no line strictly between them has an empty block list, $1\le a'<a$, and every line alternates.
\end{definition}
\cite[p.~A:31]{W13} defines normal-form proofs by their first and last lines, $\DTIME_d[n^a]$ and $\DTIME_d[n^{a'}]$ with $a'\le a$; the condition that no other line is quantifier-free is the analogue of \cite[Definition~4.2(b)]{W13}. Lemma~7.3 gives a contradiction from such a proof when $a>a'\ge1$, and Theorem~7.4 (proof sketched as ``closely follows'' Theorems~4.4 and~6.6) turns any alternation-trading proof $A_1\subseteq A_2$ between complementary simple classes into a normal-form proof. Our constructions are strict ($a'<a$). Our ceilings exclude $a'\le a$, and do not need the condition on intermediate lines (Theorem~\ref{thm:ceiling}).

\subsection{Modelling and reading choices}\label{sec:reading}
The following choices are ours or are readings of \cite{W13} on which a reader may differ.
\begin{enumerate}
\item \emph{$n^{o(1)}$ and polylogarithmic factors are dropped}; a line records exponents only, and a $\log n$ block has guess exponent $0$. This is the convention of \cite[\S4, p.~A:9]{W13} (``Polylogarithmic and $n^{o(1)}$ factors do not affect any of the elements in these proofs'').
\item \emph{Merging.} When the new block of Lemma~7.1, or the first new block of Lemma~3.3, has the same quantifier as the innermost block, the two are combined by \cite[\S2]{W13}: $(\exists n^a)^a(\exists n^b)^b=(\exists n^{\max(a,b)})^b$, so the guess exponent is the maximum and the output superscript is the new block's. Williams's proofs do this throughout; the sketch of Theorem~7.9 says ``the first 1 in a run \dots does not add a new quantifier if the class already has quantifiers'' (p.~A:37). A new block of the opposite type is \lean{convNew}; \lean{HY.Step} has no unmerged same-type block (the relaxed calculus does not track quantifiers, so it covers that case).
\item \emph{Lemma~3.3 inside a prefix.} Lemma~3.3 is stated for a quantifier-free $\DTISP$ class with outputs $\max\{1,x+e\}$ and $\max\{1,e\}$; inside a prefix whose kernel has input superscript $b$ we use $\max\{b,x+e\}$ and $\max\{b,e\}$, as Speedup Rule~1 of \cite[Definition~4.1]{W13} does for $\mathsf{DTS}$; Williams's printed 25-line proof replays under this reading (Section~\ref{sec:calibration}). The side conditions are those printed in Lemma~3.3: $a\ge1$, $0\le e\le a-x$, $0\le x\le a-1$. We call this the \emph{literal} reading. Since the kernel reads an input of length $n^{b}$, one may also require that its time stays at least its input length, i.e.\ kernel time $\ge\pin$ after every step; we call this, together with space $\le$ time on every $\DTISP$ line (item~5), the \emph{strict} reading. Theorem~\ref{thm:lower} holds under both.
\item \emph{Speedup Rule~2} (two new blocks) is included as \lean{speedNew}, and Lemma~7.1 may open a block of the opposite type anywhere (\lean{convNew}). \cite[p.~A:31]{W13} says Rule~2 can be shown unnecessary. None of our proofs uses either move except \lean{convNew} on the first line; the ceilings allow both.
\item \emph{Lemma~7.2.} TOCT states it ``for $a_i,b_i\ge1$, and all $e\ge0$''. We allow guess exponent $a_k\ge0$, because Williams's proofs slow down $\log n$ blocks (e.g.\ lines~3--4 of the 25-line proof). This changes no value, since $b_{k+1}\ge1$ is in the same maximum. The long version \cite[Lemma~5.2]{Wfull} requires $e\le a_{k+1}$ for a $\DTISP$ kernel; \lean{HY.Step} does not, and the strict-reading proofs satisfy it on every line.
\item \emph{The output of Lemma~7.2 is always $\DTIME_d$.} The sketch of Theorem~7.9 says ``After a slowdown, the deterministic portion of a class is $\DTIME_d$'' (p.~A:36). The prose proof of Theorem~7.5 writes $(\exists n^{x+d(a-1)})(\forall n)\DTISP[n^c,n^{d/(d+1)}]$ after a slowdown (p.~A:34), and its next step uses this line as time $n^c$; we read it as $\DTIME_d[n^c]$.
\item \emph{Lemma~3.4 for $\DTISP$ kernels.} \cite[p.~A:31]{W13} lists ``Slowdown via Lemma~3.4 when $C$ is a DTISP class''. Lemma~3.4 assumes $\NTIME[n]\subseteq\DTISP[n^c,n^e]$, which does not follow from $H_c$ in any way we know of for $e<c$. We take the slowdown of a $\DTISP$ kernel to be the second inclusion of Lemma~7.2, as Williams's proofs do. See Section~\ref{sec:coverage} for the version with $e=c$.
\item \emph{Normal form}: strict for the constructions, $a'\le a$ for the ceilings (Definition~\ref{def:nf} and the remark after it).
\item \emph{No weakening rule} in \lean{HY.Step}: every output is the exact one. The relaxed calculus bounds outputs only from below.
\item \emph{Starting exponent and length.} The starting exponent $a$ is free, as in Williams's linear programs, and so is the length of the proof. Our certificates start at large exponents (about $1.4\cdot10^{17}$ for $d=1$; Section~\ref{sec:stair}), which keeps every $\ge1$ floor slack, and unfold to very long proofs. Lemma~7.3 and Theorem~7.4 put no bound on either.
\end{enumerate}

\subsection{Calibration}\label{sec:calibration}
The script \texttt{w13\_25line.py} replays the 25-line proof printed in \cite[\S7.2]{W13} (pp.~A:32--A:33) under the rules of \lean{HY.Step}, at $c=1.297844$ (the value implied by its lines~9 and~10). It infers each rule from the shape of the next line and its parameter from the printed numbers, carries the output superscripts itself (they are not printed), and compares every printed number. Every step is a rule of the formal calculus, the largest deviation from a printed number is $1.0\cdot10^{-8}$ (the printed numbers have 9--10 significant digits), and the computed last line exceeds the first by $1.0\cdot10^{-8}$, so the printed proof is tight at $c=1.297844$.

\section{Statements}\label{sec:statements}

\begin{theorem}[Proofs above $r_d$; \lean{conjecture\_fails}, \lean{HStair.stairD\{1,2,3\}\_nf} and \lean{\_snf}]\label{thm:lower}
For $(d,c)=(1,\,130921174/10^8)$, $(2,\,11915968/10^7)$ and $(3,\,11386523/10^7)$ there is a strict normal-form proof \lean{HY.NF c d M n} in which every line $M_i$ has kernel time $\ge\pin$ and, if its kernel is $\DTISP[n^a,n^e]$, $e\le a$. Every real root of $p_d$ is below $c$ (\lean{HY.pd1\_root\_lt}, \lean{pd2\_root\_lt}, \lean{pd3\_root\_lt}: below $1.301$, $1.1888$, $1.1373$).
\end{theorem}

\begin{theorem}[A 39-line proof; \lean{HY.cert39\_above\_r1}, \lean{HY.cert39\_time}]\label{thm:39}
There is a strict normal-form proof with 39 lines at $c=26019/20000=1.30095$, $d=1$, from $a_0=10^5$, in which every line has kernel time $\ge\pin$; every real root of $p_1$ is below $1.30095$.
\end{theorem}

\begin{theorem}[Ceilings; \lean{HLM.ceiling\_d1}, \lean{HLM.ceiling\_d2}, \lean{HLM.ceiling\_d3}]\label{thm:ceiling}
Let $(d,c_0)=(1,6547/5000)$, $(2,149/125)$ or $(3,712/625)$, i.e.\ $c_0=1.3094$, $1.192$, $1.1392$. For every real $c\ge c_0$:
\begin{enumerate}
\item there is no strict normal-form proof \lean{HY.NF c d M n};
\item there is no relaxed path \lean{HL.RPath c d M n} with $n\ge1$ from $([\,],\lean{T}\,a_0)$, $a_0>0$, to $([\,],\lean{T}\,a')$ with $a'\le a_0$;
\item every \lean{HY.Step c d} path of length $n\ge1$ from $([\,],\lean{T}\,a_0)$, $a_0>0$, to $([\,],\lean{T}\,a')$ has $a_0<a'$.
\end{enumerate}
Parts~(2) and~(3) put no condition on the intermediate lines (they may be quantifier-free or non-alternating).
\end{theorem}

\begin{theorem}[A closed-form ceiling; \lean{HL.no\_nf}]\label{thm:sqrt}
For real $d>0$ and $c\ge1$ with $c^2\ge1+1/d$ there is no strict normal-form proof \lean{HY.NF c d M n}.
\end{theorem}

\begin{corollary}[\lean{threshold\_bounds}]\label{cor:interval}
Let $c^*_d=\sup\{c:\exists M\,n,\ \lean{HY.NF c d M n}\}$. Then $1.30921174\le c^*_1\le1.3094$, $1.1915968\le c^*_2\le1.192$ and $1.1386523\le c^*_3\le1.1392$.
\end{corollary}
\lean{threshold\_bounds} states the existence at the lower values and the non-existence for every $c$ at or above the upper values; the bounds on the supremum follow. We do not prove that the set of $c$ with a proof is an interval (the exact rules have equalities and $\ge1$ floors, so monotonicity in $c$ is not immediate); for the complexity reading this does not matter, since $\NTIME[n]\not\subseteq\DTIME_d[n^{c}]$ implies the same for every smaller exponent.

\section{The constructions}\label{sec:construction}

\subsection{A 39-line proof at \texorpdfstring{$c=1.30095$}{c = 1.30095}}\label{sec:39}
The tag list \lean{HY.cert39} (word \texttt{11101101101100111011011001101101101100}) is checked by the executable checker \lean{HY.nfCheck}, whose soundness theorem \lean{HY.nfCheck\_sound} turns acceptance into an \lean{HY.NF} proof over $\mathbb R$; the extra condition time $\ge$ input is \lean{HY.timeCheck}. Both run by \lean{decide +kernel}. The proof, with exponents divided by 1000 (so $a_0=10^5$ is shown as $100$; the certificate is exact), as printed by \texttt{replay39.py}, an independent transcription of Williams's rules:
{\footnotesize
\begin{verbatim}
 0       DTIME_1[n^100.000]
 1 L7.1  (E n^56.454)^56.454 DTISP[n^100.000, n^43.546]
 2 L3.3  (E n^56.454)^56.454 (A log)^56.454 DTISP[n^87.091, n^43.546]
 3 L3.3  (E n^56.454)^56.454 (A n^58.466)^58.466 (E log)^56.454 DTISP[n^72.171, n^43.546]
 4 L7.2  (E n^56.454)^56.454 (A n^58.466)^58.466 DTIME_1[n^93.891]
 5-13    three rounds L7.1, L3.3, L7.2 at the A level: 93.891 -> 92.172 -> 87.699 -> 76.061
14 L7.2  (E n^56.454)^56.454 DTIME_1[n^98.952]
15 L7.1  (E n^56.454)^56.454 DTISP[n^98.952, n^42.498]
16 L3.3  (E n^59.082)^59.082 (A log)^56.454 DTISP[n^82.368, n^42.498]
17 L3.3  (E n^59.082)^59.082 (A n^56.454)^56.454 (E log)^56.454 DTISP[n^68.412, n^42.498]
18 L7.2  (E n^59.082)^59.082 (A n^56.454)^56.454 DTIME_1[n^89.000]
19-24   two rounds at the A level: 89.000 -> 84.681 -> 73.444
25 L7.2  (E n^59.082)^59.082 DTIME_1[n^95.547]
26-37   four rounds at the E level: 95.547 -> 94.880 -> 93.142 -> 88.623 -> 76.862
38 L7.2  DTIME_1[n^99.994]                                  99.994 < 100
\end{verbatim}
}
Lines~0--14 are the first part of Williams's pattern: the $\exists$ block, an excursion to a $\forall$ level with three simple rounds there (lines~5--13), and the return to the $\exists$ level (line~14). Lines~15--25 make a second excursion under the same $\exists$ block: the Lemma~3.3 guess of line~16 merges into that block and raises it from $56.454$ to $59.082$, at no cost in alternations. This second excursion is what the annotation family of \cite[Theorem~7.5]{W13} does not contain. Lines~26--37 are simple rounds under the raised $\exists$ block, and line~38 is the final slowdown. The final exponent is $99994.1297<10^5$. On lines~12, 23 and~36 the kernel time exceeds the input superscript by less than $10^{-6}$ (the smallest gap is $4.6\cdot10^{-20}$, on line~12).

\subsection{Level claims}\label{sec:claims}
Proofs at $c$ close to the threshold are far too long to list, so they are checked as directed acyclic graphs of reusable sub-derivations. A sub-derivation is stated as a claim about a whole family of lines, at every scale (\texttt{Dag.lean}, namespace \lean{HD}):

\begin{definition}[\lean{HD.SC}, \lean{HD.TC}]
Call a block list \emph{good} (\lean{HD.Good}) if it alternates and every block has $1\le o$ and $0\le g\le o$. \lean{HD.Der} is a \lean{HY.Step} derivation whose lines after the first have nonempty, alternating block lists. The \emph{S-claim} \lean{SC}$(c,d,B_0,\kappa,\varepsilon,V)$ says: for every scale $B\ge B_0$ and every line $(u::\mathit{rest},\DTISP[n^a,n^e])$ with good blocks, $u.o\le B$, $1\le a\le\kappa B$ and $0\le e\le\varepsilon B$, there are $a'$ with $1\le a'\le c\max(\pin\,\mathit{rest},VB)$, a \lean{Der} derivation from that line to some line $L_p$, and one \lean{HY.Step} from $L_p$ to $(\mathit{rest},\DTIME_d[n^{a'}])$. The \emph{T-claim} \lean{TC}$(c,d,B_0,\alpha,V)$ says the same for lines $(u::\mathit{rest},\DTIME_d[n^a])$ with $1\le a\le\alpha B$.
\end{definition}
A claim is about the level of the innermost block: it says that the level can be closed, removing that block, at cost $c\max(\pin\,\mathit{rest},VB)$. Because claims are downward closed (they speak about every good line whose top, time and space are at most the planned ones) and hold at every scale $B\ge B_0$, one claim can be used at every call site. Claims are combined by three lemmas, all proved in Lean for the exact rules:
\begin{itemize}\sloppy
\item \lean{HD.SC\_pop}, \lean{HD.TC\_pop}: closing the level at once (one Lemma~7.2 step) proves $\lean{SC}(\dots,\max(1,\kappa))$ and $\lean{TC}(\dots,\max(1,\alpha))$, for $c\ge1$ and $B_0\ge0$.
\item \lean{HD.TC\_of\_SC} (Lemma~7.1, merged): for $B_0\ge0$, $d>0$, $\kappa\ge1$, $\varepsilon>0$ and $d(\kappa-1)\le\varepsilon$, an S-claim $\lean{SC}(B_0,\kappa,\varepsilon,\Omega)$ gives $\lean{TC}(B_0,\alpha,\Omega\alpha/\kappa)$ for every $\alpha\ge\kappa$. The T-level is rescaled to top $\beta=\alpha B/\kappa$ by Lemma~7.1 with $s=\min(a,\varepsilon\beta/d)$ if $a\le\beta$ and $s=a-\beta$ otherwise.
\item \lean{HD.SC\_speed} (Lemma~3.3, merged): for $c\ge1$ and $B_0\ge0$, from an inner S-claim $\lean{SC}(B_0,\kappa',\varepsilon,W)$ and an outer T-claim $\lean{TC}(B_0,\alpha,V_T)$, with $1\le u$, $\varepsilon\le u$, $\kappa'=\kappa+\varepsilon-u$, $\varepsilon\le\kappa'$, $0\le\varepsilon\le1$, $1\le\kappa'B_0$, $c\max(u,W)\le\alpha u$ and $1\le uV_T$, follows $\lean{SC}(B_0,\kappa,\varepsilon,uV_T)$. The derivation applies one Lemma~3.3 step that raises the innermost block to top $uB$, then the inner claim for the new $\log n$ level at scale $B$, then the outer claim for the raised level at scale $uB$.
\end{itemize}
When the actual line is below the planned one, \lean{SC\_speed} clips the speedup to $x=\min((u-\varepsilon)B,\ a-\max(e,1))$, and closes the level at once if $e>a$. Every side condition of the exact rules, including the $\ge1$ floors, is discharged; no weakening rule is used. Finally \lean{HD.nf\_of\_SC}: for $c\ge1$, $B_0\ge0$, $d\ge0$, $\kappa>1$, $d(\kappa-1)\le\varepsilon$, $V\ge1$ and $cV<\kappa$, an S-claim $\lean{SC}(B_0,\kappa,\varepsilon,V)$ gives an \lean{HY.NF} proof: from $\DTIME_d[n^{\kappa\beta}]$, $\beta=\max(B_0,1)$, apply Lemma~7.1 (\lean{convNew}) with $s=\kappa\beta-\beta$, then the claim; the final line has exponent at most $cV\beta<\kappa\beta$.

\subsection{The checker and its soundness}\label{sec:checker}
\lean{HD.dagCheck c d M SH nodes root} is an executable function over $\mathbb Q$, with a grid $M$ and values in units of $2^{-\mathit{SH}}$. A node (\lean{HD.Node}) claims an S-claim at $\kappa=t/M$, $\varepsilon=d(t-M)/M$ (the space Lemma~7.1 leaves when it rescales a kernel $\kappa B$ to top $B$) with value $V=\varphi\kappa$, $\varphi=\mathit{phi}/2^{\mathit{SH}}$. Its \lean{steps} list a chain of Lemma~3.3 steps: for each, the new top $u=k_u/M$, the closing of the outer T-level after the inner level (\lean{r = 0}: Lemma~7.2 at once, value $\max(1,\alpha)$; otherwise the table entry of index $r-1$, an earlier node, which needs its $\kappa_{r-1}\le\alpha$ and gives value $\varphi_{r-1}\alpha$; function \lean{tval}), and a claimed value $w/2^{\mathit{SH}}$. \lean{chain} walks the chain: each step moves the kernel position from $j$ to $j+d(t-M)-k_u$ (which must stay $\ge1$), the innermost level is closed at the recorded final position $p$ with value $\max(1,p/M)$, and every claimed value must satisfy $1\le u\,v\le w/2^{\mathit{SH}}$ with $v$ the closing value of the outer T-level at $\alpha=c\max(1,W/u)$, $W$ the value of the rest of the chain. \lean{procNodes} checks the nodes in order ($M<t$, $d(t-M)\le M$, chain value$/\kappa\le\varphi$, $\varphi\kappa\ge1$) and records the proved claims (invariant \lean{TabOK}). \lean{dagCheck} also requires $1\le c$, $1\le d$, $1\le M$ and $c\varphi_{\mathrm{root}}<1$.

\begin{theorem}[\lean{HD.dagCheck\_sound}]\label{thm:dag}
If \lean{HD.dagCheck c d M SH nodes root = true} then there is a strict normal-form proof \lean{HY.NF c d Mp n} (with $c$ and $d$ cast to $\mathbb R$).
\end{theorem}
The proof composes \lean{tval\_sound}, \lean{chain\_sound} and \lean{procNodes\_sound}, which interpret each check by the lemmas of Section~\ref{sec:claims} at scale $B_0=M$, and ends with \lean{nf\_of\_SC}. So the normal-form proof starts at $a_0=\kappa_{\mathrm{root}}M=t_{\mathrm{root}}$.

\paragraph{The strict reading} (\texttt{DagStrict.lean}, namespace \lean{HDS}). The same claims and checker with: \lean{HDS.Der} requires kernel time $\ge\pin$ and \lean{SpOK} (space $\le$ time on a $\DTISP$ kernel) on every line after the first; claims assume $u.o\le a$ and $e\le a$ on their first line and guarantee $\pin\,\mathit{rest}\le a'$; a speedup is clipped at $x\le a-\max(e,u.o,1)$; chain positions must stay $\ge M$ ($\kappa'\ge1$); Lemma~7.1 picks $s\le a/d$. \lean{HDS.dagCheck\_sound} gives an \lean{HY.NF} proof with kernel time $\ge\pin$ and \lean{SpOK} on every line.

\subsection{The staircase certificates}\label{sec:stair}
The three certificates of Theorem~\ref{thm:lower} (\texttt{StairD1.lean}, \texttt{StairD2.lean}, \texttt{StairD3.lean}) are node lists produced by an untrusted generator (\texttt{scripts/stair\_cert.py}) from the level structure of Section~\ref{sec:threshold}. Each list is checked twice by \lean{decide +kernel}, by \lean{HD.dagCheck} (\lean{stairD\{d\}\_ok}) and by \lean{HDS.dagCheck} (\lean{stairD\{d\}\_sok}). Lean checks that the node lists are proofs; it does not check that they are the proofs the recursion describes.

\begin{table}[ht]
\centering
\begin{tabular}{cccccccc}
\toprule
$d$ & $c$ & grid $M$ & units & nodes & speedup steps & root $c\varphi$ & unfolded lines \\
\midrule
1 & $130921174/10^8$ & $2^{56}$ & $2^{-120}$ & 552 & 3915 & 0.99916308 & $\le7.61\cdot10^{21}$ \\
2 & $11915968/10^7$ & $2^{50}$ & $2^{-110}$ & 286 & 1846 & 0.99940456 & $\le2.46\cdot10^{17}$ \\
3 & $11386523/10^7$ & $2^{50}$ & $2^{-110}$ & 212 & 1255 & 0.99961709 & $\le1.10\cdot10^{15}$ \\
\bottomrule
\end{tabular}
\caption{The staircase certificates (\texttt{scripts/dag\_check.py}, which re-implements both checkers in exact rationals on the Lean data). ``Unfolded lines'' counts the lines of the derivation that the soundness proof builds, if no level is closed early (an upper bound). The proofs start at $a_0=t_{\mathrm{root}}\approx1.41\cdot10^{17}$, $1.68\cdot10^{15}$, $1.50\cdot10^{15}$.}
\label{tab:certs}
\end{table}

The structure of the certificates (as built by the generator). In a node's chain the first step is the outermost one, so its T-level is closed last. Level~0 is a climb of nodes, each one Lemma~3.3 step with $u\approx1$ (top kept) whose inner level is closed at once and whose T-level is closed by the previous node (the first node closes it at once). These are Williams's simple rounds, which contract the kernel below the repelling fixed point $\alpha_1=c(1+d)/(c(1+d)-1)$. Level $h+1$ starts with a \emph{fresh} node, whose chain closes its T-levels, from the outermost in, with the top node (the \emph{kink}) of level $h$, again with that kink, and then with the kinks of levels $h-1,\dots,0$. It continues with a \emph{climb}, each node of which closes its outermost T-level with the previous node and the others with the kinks of levels $h,\dots,0$. The construction stops at the first kink with $c\varphi<1$; that node is the root.

\section{The ceilings}\label{sec:ceiling}

\subsection{The relaxed calculus}\label{sec:relaxed}
The ceilings are proved for a relaxation \lean{HL.RStep c d} (\texttt{Relaxed.lean}). A relaxed line keeps only the output superscripts of its blocks (innermost first) and the kernel; quantifiers and guess exponents are dropped, the input superscript of the empty prefix is $0$, and every output is bounded only from below:
\begin{itemize}
\item \lean{convM}: $(b::r,\lean{T}\,a)\to(b_1::r,\lean{S}\,a_2\,e)$ with $0\le s\le a$, $\max(b,a-s)\le b_1$, $a\le a_2$, $ds\le e$; \lean{convN}: $(bl,\lean{T}\,a)\to(g::bl,\lean{S}\,a_2\,e)$ with $\max(\pin\,bl,a-s)\le g$ and the same other conditions;
\item \lean{speedM}: $(b::r,\lean{S}\,a\,e)\to(L::b'::r,\lean{S}\,a_2\,e_2)$ with $0\le x$, $e\le a-x$, $\max(b,x+e)\le b'$, $\max(b,e)\le L$, $a-x\le a_2$, $e\le e_2$; \lean{speedN}: the same with two new blocks $g_2::g_1::bl$;
\item \lean{slow}: $(b::r,k)\to(r,\lean{T}\,a')$ with $c\max(\pin\,r,b,\mathrm{time}(k))\le a'$.
\end{itemize}
\begin{lemma}[\lean{HL.sim}]\label{lem:sim}
For $c\ge0$, every \lean{HY.Step c d L L'} gives \lean{HL.RStep c d (enc L) (enc L')}, where \lean{enc} keeps the output superscripts and the kernel.
\end{lemma}
The relaxation drops the side conditions $a\ge1$, $x\le a-1$, $s>0$, all $\ge1$ floors and the guess exponent in the slowdown maximum; since the relaxed slowdown only bounds its output from below, dropping a term only weakens its condition.

\subsection{A one-level potential: the ceiling \texorpdfstring{$c^2\ge1+1/d$}{c² ≥ 1 + 1/d}}\label{sec:onelevel}
Theorem~\ref{thm:sqrt} has a short proof that shows the method. Put $A=\max(1+1/d,c)$ and $g=1/A$. On a level whose innermost block has superscript $b$, let
\[\psi_T(b,a)=\min\bigl(\max(b,a),\ c\max(b,ga)\bigr),\qquad \psi_S(b,a,e)=\min\bigl(\max(b,a),\ c\max(b,e)\bigr)\]
for kernels $\DTIME_d[n^a]$ and $\DTISP[n^a,n^e]$. The potential $\Phi$ of a line (\lean{HL.Phi}) lifts the innermost value $y$ through each outer block of superscript $b'$ by $y\mapsto\psi_T(b',c\max(b',y))$, and finally by $y\mapsto c\max(0,y)$; a quantifier-free line $\DTIME_d[n^a]$ has potential $a$. Intuitively $\psi$ is a lower bound for what the level can still reach: an S-level can at best end at $c\max(\text{top},\text{space})$, a T-level with kernel $a$ at best at $c\max(\text{top},a/A)$. The facts used (\lean{HL.Facts}) are $1\le c$, $cg\le1$, $c^2g\ge1$ and $g\le d(1-g)$, i.e.\ $c\le A\le c^2$ and $A\ge1+1/d$; with $c^2\ge1+1/d$ they hold.

\begin{lemma}[\lean{HL.phi\_step}]
Under \lean{HL.Facts c d g}, an \lean{RStep c d} from a line with nonnegative entries does not lower $\Phi$.
\end{lemma}
\begin{proof}[Proof sketch]
Each constructor reduces to a local inequality; the lift is monotone (\lean{up\_mono}).
\emph{Lemma~7.1} (\lean{conv\_ineq}): $\psi_T(b,a)\le\psi_S(\max(b,a-s),a,ds)$ for $0\le s\le a$, because $ga\le\max(a-s,ds)$: if $s\le(1-g)a$ then $a-s\ge ga$, and otherwise $ds>d(1-g)a\ge ga$.
\emph{Lemma~3.3} (\lean{speed\_ineq}): with $b'=\max(b,x+e)$ and $L=\max(b,e)$, the inner value is $S=\min(\max(L,a-x),cL)$, and the lifted value is at least $c\max(b',g\,c\max(b',S))$. If $S=cL$, this is $\ge c^2g\cdot cL\ge cL\ge\psi_S(b,a,e)$. Otherwise $S=\max(L,a-x)$; the cases $a\le b$ and $e\le b$ are immediate, and for $a,e\ge b$ the lifted value is at least $c\,M'$ with $M'=\max(x+e,\ cg(a-x))$. The guess $x+e$ and the remaining time $a-x$ trade off: if $e\le cga$ then $(1+cg)M'\ge cg(x+e)+cg(a-x)=cg(a+e)\ge(1+cg)e$, and otherwise $M'\ge x+e\ge e$. Either way $cM'\ge ce\ge\psi_S(b,a,e)$.
\emph{Opening a block} (\lean{push\_ineq}, for \lean{convN} and \lean{speedN} under a prefix): $y\le\psi_T(b,c\max(b,y))$ for $y\ge b$, by $c^2g\ge1$.
\emph{Lemma~7.2}: $\psi\le\max(\text{top},\text{time})$ (\lean{pot\_le}), so the new exponent bounds the lifted value.
\emph{First step} (\lean{first\_ineq}): from $\DTIME_d[n^{a}]$ the potential becomes at least $c\,\psi_T(0,a)=c\min(a,cga)\ge a$.
\end{proof}
So along every relaxed path from $\DTIME_d[n^{a_0}]$, $a_0\ge0$, every later quantifier-free line $\DTIME_d[n^{a'}]$ has $a'\ge a_0$ (\lean{HL.qf\_ge}), and with Lemma~\ref{lem:sim} no strict normal-form proof exists (\lean{HL.no\_nf}). The bound $\sqrt{1+1/d}$ is the previous best exponent $\sqrt{(d+2)/(d+1)}$ recalled in \cite[p.~A:3]{W13} (for $d=1$ the van Melkebeek--Raz bound $\sqrt{3/2}$, p.~A:31) with $d$ shifted by one.

\subsection{Multi-level potentials}\label{sec:multilevel}
The one-level potential ignores that a T-level can be closed more cheaply by running further subproofs under it. The multi-level potential (\texttt{Potential.lean}, namespace \lean{HLM}) has a finite list of levels $(P_j,\theta_j)$ together with level~0, $(0,\theta_0)$; put $\kappa_p=\max(c,P_p)$. Then
\begin{align*}
\psi_T(b,a)&=\min\bigl(\max(b,a),\ \max(cb,\ \mathrm{Tm}(b,a))\bigr),\qquad \mathrm{Tm}(b,a)=\min_j\max(P_jb,\ \theta_ja/c),\\
\psi_S(b,a,e)&=\inf_w\max\bigl(0,V_w(b,a,e)\bigr),
\end{align*}
where $w$ ranges over words (lists of levels, outermost first), $V_{[\,]}(b,a,e)=\max(b,a)$ and
\[V_{p::w}(b,a,e)=\max\Bigl(\kappa_pb,\ \theta_p\,\mathrm{Ad}(w)\max(b,e),\ \frac{a+|p::w|\,e}{D(p::w)}\Bigr),\]
with $\mathrm{Ad}([\,])=1$, $\mathrm{Ad}(p::w)=\max(\kappa_p,\theta_p\mathrm{Ad}(w))$, $D([\,])=1$, $D(p::w)=1/\kappa_p+D(w)/\theta_p$. A word stands for a chain of Lemma~3.3 steps whose T-levels are closed by the listed levels; $V_w$ is an explicit lower bound for the cost of that chain. The one-level potential of Section~\ref{sec:onelevel} has a cruder $\psi_S$; with only level~0 and $\theta_0=c^2/A$ the T-parts agree.

\lean{HLM.Base}: $1\le c$, $0\le d$, and $1\le\theta_p\le c$, $P_p\ge0$ for every level. \lean{HLM.Hyp} adds, for every nonempty word $w$, the condition \lean{ConvOK} for its state $(\mathrm{Ad}\,w,\mathrm{AL}\,w,D\,w,|w|)$, where $\mathrm{AL}([\,])=0$, $\mathrm{AL}(p::w)=\theta_p\mathrm{Ad}(w)$, and $\lean{ConvOK}(A,L,D,m)$ says $\mathrm{Tm}(b,a)\le\max(Ab,\ Q\,a)$ for all $b,a\ge0$, with
\[Q=\max\Bigl(\frac{ALd}{A+Ld},\ \frac{A(1+md)}{AD+md}\Bigr).\]

\begin{lemma}[\lean{HLM.phi\_step}]
Under \lean{HLM.Hyp}$(c,d,\theta_0,\mathit{lv})$, an \lean{HL.RStep c d} from a line with nonnegative entries does not lower $\Phi$.
\end{lemma}
\begin{proof}[Proof sketch]
\emph{Lemma~3.3} needs only \lean{Base}. Prepending a letter absorbs one speedup (\lean{V\_succ}): for $b,x\ge0$,
\[V_{p::w}(b,a,e)\le\max\bigl(\kappa_p\max(b,x+e),\ \theta_pV_w(\max(b,e),a-x,e)\bigr).\]
The first two terms of $V_{p::w}$ are bounded directly ($V_w(L,\cdot)\ge\mathrm{Ad}(w)L$), and the third by a weighted average with weights $1/\kappa_p$ and $D(w)/\theta_p$: $a+(m+1)e=(x+e)+(a-x+me)$ (\lean{U\_succ}). The lift satisfies $\psi_T(b',c\max(b',W))\ge\min_p\max(\kappa_pb',\theta_pW)$ (\lean{le\_psT\_lift}); together they give \lean{speed\_ineq}.
\emph{Lemma~7.1} (\lean{conv\_word}, \lean{conv\_ineq}): for a nonempty word, $V_w(\max(b,a-s),a,ds)\ge\max(\mathrm{Ad}\cdot\max(b,a-s),\ \mathrm{AL}\cdot ds,\ (a+|w|ds)/D)$, and two weighted averages over $s$ (\lean{cross}) give $\ge\max(\mathrm{Ad}\,b,\ Q\,a)$, which dominates $\psi_T(b,a)$ by \lean{ConvOK} and $\mathrm{Ad}\ge c$. The empty word is $\max(b,a)\ge\psi_T$.
\emph{Opening a block} (\lean{push\_ineq}) uses $\theta\ge1$; \emph{Lemma~7.2} and the \emph{first step} (\lean{first\_ineq}, using $\mathrm{Tm}(0,a)\ge a/c$) are as before.
\end{proof}
Consequences: \lean{HLM.qf\_ge} ($a_0\le a'$ along relaxed paths from $\DTIME_d[n^{a_0}]$, $a_0\ge0$), \lean{HLM.no\_nf}, and, when every $\theta$ (level~0 included) exceeds $1$ and $c>1$, \lean{HLM.qf\_gt} and \lean{HLM.no\_nf\_le}: the first step raises the potential strictly (\lean{first\_strict}), so $a_0<a'$ along every path of length $\ge1$ from $a_0>0$, which excludes non-strict normal forms.

\paragraph{All larger $c$.} The relaxed calculus is monotone in $c$ (\lean{HLM.rstep\_mono}): an \lean{RStep} at $c$ from a line with nonnegative entries is one at every $c_0\le c$, because only \lean{slow} depends on $c$ and its lower bound only weakens. So a potential valid at $c_0$ excludes proofs at every $c\ge c_0$ (\lean{rge\_of}, \lean{rgt\_of}, \lean{no\_nf\_of}, \lean{no\_nf\_le\_of}). For the exact calculus we neither claim nor use it; it enters only through \lean{HL.sim}.

\subsection{Checking the Lemma 7.1 condition for all words}\label{sec:abs}
\lean{Hyp} quantifies over infinitely many words. \texttt{AbsCheck.lean} reduces it to a finite check. Prepending a letter $p$ maps a state $(A,L,D,m)$ to $(\max(\kappa_p,\theta_pA),\ \theta_pA,\ 1/\kappa_p+D/\theta_p,\ m+1)$. This map is monotone for the order ``$A$, $L$, $m$ larger and $D$ smaller'' (\lean{Dom}, \lean{dom\_prep}), and \lean{ConvOK} transfers along that order (\lean{convOK\_mono}; monotonicity of $Q$ in $m$ uses $AD\ge1$, which holds for every real word, \lean{AdDw\_ge\_one}). A state is \emph{done} once $\theta_0/c\le ALd/(A+Ld)$: then \lean{ConvOK} holds (\lean{convOK\_done}) and stays true under prepending. The executable \lean{absCheck}$(c,d,\theta_0,\mathit{lv},S)$ accepts when $1\le c$, $0\le d$, every level has $1\le\theta\le c$ and $P\ge0$, every state of the list $S$ has $A>0$, $L\ge0$, $D>0$ and an interval certificate for \lean{ConvOK} (\lean{ivOK}: breakpoints in the ratio $a/b$, starting at $0$, with one level per piece whose term is below $\max(Ab,Qa)$ on that piece), and $S$ together with the empty word's state $(1,0,1,0)$ is closed under prepending each letter, up to done states and domination of a state of $S$.

\begin{theorem}[\lean{HLM.absCheck\_sound}]
If \lean{absCheck}$(c,d,\theta_0,\mathit{lv},S)={}$\lean{true} then \lean{HLM.Hyp} holds for the data cast to $\mathbb R$.
\end{theorem}
The proof is an induction over the real word (prepending), so word length is unbounded.

\paragraph{The three ceilings} (\texttt{Ceiling.lean}). In each case $\theta_0=c_0^2-c_0/(d+1)$; the further levels have $P_{j+1}=\theta_jP_j$ with $P_1=c_0\theta_0$, and $\theta_j$ slightly above $1$, found by an untrusted search. Only their acceptance matters.
\begin{center}
\begin{tabular}{cccl}
\toprule
$d$ & $c_0$ & abstract states & levels $(P_j,\theta_j)$, $j\ge1$ \\
\midrule
1 & $6547/5000$ & 105 & $(1.38774,1.027)$, $(1.42521,1.017)$, $(1.44944,1.014)$ \\
2 & $149/125$ & 50 & $(1.22005,1.011)$, $(1.23347,1.008)$ \\
3 & $712/625$ & 19 & $(1.15398,1.0055)$ \\
\bottomrule
\end{tabular}
\end{center}
Each list is accepted by \lean{decide +kernel} (\lean{absD\{d\}\_ok}); every $\theta$ exceeds $1$ (\lean{thD\{d\}\_gt}); and \lean{ceiling\_d\{d\}} combines \lean{no\_nf\_of}, \lean{no\_nf\_le\_of} and \lean{rgt\_of} into Theorem~\ref{thm:ceiling}. The script \texttt{abs\_check.py} re-implements \lean{absCheck} in exact rationals on the Lean data; it accepts all three, and rejects each when $c_0$ is lowered by $10^{-4}$ (with $\theta_0$ recomputed) or when the last abstract state is removed.

\section{What the relaxed calculus covers}\label{sec:coverage}
Lemma~\ref{lem:sim} embeds \lean{HY.Step}. Beyond it, \lean{HL.RStep} contains:
\begin{itemize}
\item weakening of every output, time and space exponent (all are bounded only from below);
\item Lemma~7.1 with $s=0$, and Lemma~3.3 without $a\ge1$, $x\le a-1$ and $e\ge0$;
\item Speedup Rule~2 and Lemma~7.1 opening a block anywhere (\lean{speedN}, \lean{convN}), and since quantifiers are not tracked, new blocks of either type, merged or not;
\item Lemma~3.3 on a quantifier-free $\DTISP$ kernel (Speedup Rule~0), which cannot arise from a $\DTIME_d$ line anyway;
\item no $\ge1$ floors, input superscript $0$ for the empty prefix, and the slowdown without the guess exponent;
\item paths whose intermediate lines are quantifier-free or do not alternate (Theorem~\ref{thm:ceiling}(2), (3)).
\end{itemize}
Combining adjacent same-type blocks after the fact is covered too (a paper argument): all rules act on the innermost block and never change outer blocks, so a same-type adjacency arises only when a block is created, and ``create unmerged, then combine'' gives the merged rule's result (the combined block's output is the inner one's).

Not covered:
\begin{itemize}
\item \emph{The inclusion $\DTIME_d[t]\subseteq\DTISP[t,t]$ inside a prefix}, for $d\ge2$. It is sound (a $d$-dimensional tape machine running $t$ steps visits at most $t$ cells), but it is not among Williams's rules (p.~A:31 lists only Lemma~3.3 or~7.1 for speedups and Lemma~3.4 or~7.2 for slowdowns). Lemma~3.4 with $e=c$, the one form of Lemma~3.4 that $H_c$ supports (item~7 of Section~\ref{sec:reading}), gives the same result as Lemma~7.2 followed by this inclusion. Inside a prefix and for $d=1$ it is the case $s=a$ of \lean{convM}, so the ceiling for $d=1$ covers it there; on a quantifier-free line, and for $d=2,3$ everywhere, it is outside \lean{RStep}, and Theorem~\ref{thm:ceiling} does not cover proofs that use it. The one-level argument of Section~\ref{sec:onelevel} extends to it on paper, since $\psi_S(b,a,a)=\max(b,a)\ge\psi_T(b,a)$; the multi-level ceilings have not been checked against it.
\item Other sound moves that are not rules of the calculus, e.g.\ $\DTISP[t,s]\subseteq\DTIME_d[t\cdot s^{1/d}]$ (laying out the space as a $d$-cube), other speedup or slowdown lemmas, and contradictions other than a normal-form proof. Extending the ceilings to arbitrary alternation-trading proofs between complementary classes uses \cite[Theorem~7.4]{W13}, whose proof is only sketched there.
\end{itemize}

\section{The threshold}\label{sec:threshold}

\subsection{What is proved}
By Corollary~\ref{cor:interval}, $c^*_1\in[1.30921174,1.3094]$ (width $1.9\cdot10^{-4}$), $c^*_2\in[1.1915968,1.192]$ ($4.0\cdot10^{-4}$), $c^*_3\in[1.1386523,1.1392]$ ($5.5\cdot10^{-4}$). The exact values are open.

\subsection{The staircase recursion (numerics, not proved)}\label{sec:recursion}
Fix a level with top $B$ and normalise $B=1$. Numerical optimisation over the reduced system of level values suggests that the least achievable closing value $F(\alpha)$ of a T-level with kernel $\alpha$ is a staircase of constant pieces: level~0 is Williams's simple-round basin, with value $\Psi_0=c$ up to the kink $r_0=\alpha_1=c(1+d)/(c(1+d)-1)$, and each further level $h$ has a value $\Psi_h$ on a branch ending at a kink $r_h$, with ratio $g_h=\Psi_h/r_h$. A point of a branch is reached by one Lemma~3.3 step with $u=1$ into a lower point of the same branch, then one step into each lower kink; a new level starts with the fresh chain of Section~\ref{sec:stair}. This leads to the recursion (\texttt{staircase.py}): $S_{-1}=1$, $S_0=S_{-1}/(cg_0)+1/\Psi_0$, and for $h\ge0$, with $n_h=1+(h+2)d$,
\[\Psi_{h+1}=c\,g_h\Psi_h,\qquad r_{h+1}=\max\Bigl(\frac{n_h}{n_h-S_h/c},\ \frac{cg_h+n_h+(r_h/c)S_{h-1}}{n_h}\Bigr),\qquad g_{h+1}=\frac{\Psi_{h+1}}{r_{h+1}},\]
$S_{h+1}=S_h/(cg_{h+1})+1/\Psi_{h+1}$. The first term of $r_{h+1}$ is the kink of the branch, the second the end of the fresh chain's kernel budget. Level $h+1$ exists iff $n_h>S_h/c$, $r_{h+1}\le1+1/d$ and $r_{h+1}>c\Psi_h$ (equivalently $g_{h+1}<g_h$). The recursion \emph{refutes} at $c$ if some existing level has $cg_h<1$, and \emph{stops} if the staircase ends with $cg_h>1$ throughout. At $c=1.3092$, $d=1$, it gives $r_h=1.61789$, $1.76931$, $1.83631$, $1.87385$, $1.89816$, $1.91635$, $1.93374$ and $cg_h=1.0594$, $1.0263$, $1.0148$, $1.0093$, $1.0056$, $1.0016$, $0.9941$, so it refutes at level~6. The generator of Section~\ref{sec:stair} builds certificates from exactly this structure, and they check; that is the only machine-checked link between the recursion and the calculus.

\begin{definition}
$\cstair(d)$ is the boundary between the $c$ at which the recursion refutes and those at which it stops.
\end{definition}
\texttt{staircase.py threshold} bisects in decimal interval arithmetic with directed rounding, and certifies the verdict at both ends of the final bracket (that the verdict is monotone in $c$ is observed, not proved):
\begin{align*}
\cstair(1)&=1.3092117446776010510221902603497164170365731409176335370894\ldots\\
\cstair(2)&=1.19159683545302011311694956315014944535226780478297581606997\ldots\\
\cstair(3)&=1.13865239000087119328135666360349492168754860133162470872070\ldots
\end{align*}
(the digits shared by both ends of a bracket of width $10^{-60}$; 116--149 levels are needed at the bracket ends, and near the threshold the first term of $r_{h+1}$ is the maximum). The $d=1$ certificate of Theorem~\ref{thm:lower} is $4.7\cdot10^{-9}$ below $\cstair(1)$. Numerically, the orbit is of saddle type near the threshold: below it $cg_h$ eventually drops below $1$; above it $g_h$ eventually increases and the staircase stops. We see no reason for $\cstair$ to be algebraic and did not search for a closed form.

Level $h$ exists for $c$ up to $c_h$ (\texttt{staircase.py appear}); for $d=1$, $c_1,\dots,c_6=1.330278$, $1.311714$, $1.309609$, $1.309287$, $1.309228$, $1.309215$, decreasing to $\cstair(1)$. The first value, $c_1=1.330277691377$, $1.199453639736$, $1.142696215782$ for $d=1,2,3$, coincides to the 12 printed digits with the real root in $[1.1,2]$ of $8c^3-2c^2-13c+2$, $36c^3-9c^2-46c+6$ and $96c^3-20c^2-113c+12$ respectively (\texttt{staircase.py roots}); we have no proof of this identity. Heuristically, a potential with levels $0,\dots,h$ (the level values of the staircase) should be valid as long as level $h+1$ does not exist, i.e.\ for $c\ge c_{h+1}$. The checked ceilings are consistent with this and lag a little behind: the four-level ceiling for $d=1$ is $1.3094$, above $c_4=1.309287$; the three-level ceiling for $d=2$ is $1.192$, above $c_3=1.191768$; the two-level ceiling for $d=3$ is $1.1392$, above $c_2=1.139159$.

\begin{conjecture}\label{conj}
$c^*_d=\cstair(d)$ for $d=1,2,3$.
\end{conjecture}
\emph{Evidence}: the certificates built from the recursion check at $c$ within $4.7\cdot10^{-9}$ of $\cstair(1)$ (and at $c$ within $3.5\cdot10^{-8}$ and $9.0\cdot10^{-8}$ of $\cstair(2)$ and $\cstair(3)$); and the ceilings, which mirror the levels of the staircase, come down toward it as levels are added ($\sqrt2$ with one level, $1.3094$ with four, for $d=1$). \emph{Caveats}: the recursion is a derivation from numerically optimal proofs, not a theorem; we have not proved that a refuting verdict always yields a proof, or that the recursion contains every optimal proof shape; and the ceilings are equally consistent with any value of $c^*_d$ in the interval of Corollary~\ref{cor:interval}. The conjecture should be treated as open.

\section{Consequence for lower bounds}\label{sec:consequence}

\begin{corollary}[conditional; paper argument]\label{cor:lb}
Assume the following results of \cite{W13} for the hybrid model with a $d$-dimensional tape: Lemma~7.1 (for both quantifiers), Lemma~3.3 applied inside a prefix in the form of Speedup Rule~1, Lemma~7.2 (also for $\log n$ blocks), Lemma~7.3, the combining of same-type blocks of \S2, and Theorem~3.1 with Corollary~3.2. Then $\NTIME[n]\not\subseteq\DTIME_d[n^{c}]$ for $c=1.30921174$, $1.1915968$, $1.1386523$ ($d=1,2,3$), and hence $\SAT\notin\DTIME_d[n^{c}]$ for every $c<1.30921174$, $1.1915968$, $1.1386523$ respectively. In particular $\SAT\notin\DTIME_1[n^{1.3092}]$, $\SAT\notin\DTIME_2[n^{1.1915}]$ and $\SAT\notin\DTIME_3[n^{1.1386}]$.
\end{corollary}
\begin{proof}[Argument]
Under $H_c$, each line of a normal-form proof of Theorem~\ref{thm:lower} is contained in the next by the named lemma (Table~\ref{tab:rules}, with the readings of Section~\ref{sec:reading}), so $\DTIME_d[n^{a}]\subseteq\DTIME_d[n^{a'}]$ with $a>a'\ge1$, contradicting Lemma~7.3. The proof has a fixed finite length, so the $n^{o(1)}$ losses of its steps compose to $n^{o(1)}$, and $a-a'$ is a fixed positive constant. Hence $\NTIME[n]\not\subseteq\DTIME_d[n^c]$, which implies the same for every smaller exponent, and Corollary~3.2 of \cite{W13} transfers it to SAT up to a polylogarithmic factor.
\end{proof}
This is exactly the dependence of \cite[Theorem~7.5]{W13}. Lemma~7.1 is proved in \cite{W13} as a sketch, for $d=1$ (the universal case in one sentence) with a further sketch for $d\ge2$ (pp.~A:29--A:30); Lemma~7.2's proof ``directly follows that of earlier, similar results'' (p.~A:31); Lemma~7.3 is ``analogous'' to Lemmas~4.3 and~6.3. The strict-reading proofs also satisfy the condition $e\le a_{k+1}$ of \cite[Lemma~5.2]{Wfull} and kernel time $\ge$ input length on every line. Our proofs use the lemmas only in configurations that occur in Williams's printed proofs, but many more times (constant in $n$, though astronomically large). The normal-form theorem, Theorem~7.4, is not needed for the corollary, because the proofs are already in normal form; it is needed to read the ceilings as statements about all alternation-trading proofs.

\section{What is and is not verified}\label{sec:verified}

\paragraph{(a) Kernel-checked.}
Lean~4.31.0 with Mathlib (pinned revision). The library builds with no \texttt{sorry}, no \texttt{admit}, no \texttt{native\_decide} and no new axioms; \texttt{\#print axioms} reports only \texttt{propext}, \texttt{Classical.choice} and \texttt{Quot.sound} for every headline theorem (\texttt{Axioms.lean} in \texttt{PastaMath/\allowbreak OneTapeThreshold/}). The certificates are checked with \lean{decide +kernel}, which adds no axiom. Checked are: Theorems~\ref{thm:lower}, \ref{thm:39}, \ref{thm:ceiling}, \ref{thm:sqrt} and Corollary~\ref{cor:interval} for the formal calculus; the soundness of the three checkers (\lean{HY.nfCheck\_sound}, \lean{HD.dagCheck\_sound}, \lean{HDS.dagCheck\_sound}) and of the abstraction (\lean{HLM.absCheck\_sound}); the simulation \lean{HL.sim}; monotonicity of the potentials and of the relaxed calculus in $c$; and the root bounds for $p_d$.

\paragraph{(b) Not machine-checked.}
\begin{itemize}
\item \emph{That the formal rules are Williams's.} This is a modelling claim, compared rule by rule in Table~\ref{tab:rules}, with the reading choices of Section~\ref{sec:reading} and the calibration of Section~\ref{sec:calibration}. The constructions use only \lean{HY.Step}; the ceilings are proved for the relaxation of Section~\ref{sec:relaxed}, with the gaps listed in Section~\ref{sec:coverage}.
\item \emph{Soundness of Williams's lemmas} (Section~\ref{sec:consequence}) and Corollary~\ref{cor:lb}. We do not formalise complexity classes.
\item \emph{The staircase recursion, $\cstair$ and Conjecture~\ref{conj}} (Section~\ref{sec:recursion}): numerics. The 58--59-digit values are certified for the recursion's verdicts by interval arithmetic, not for the calculus.
\item \emph{The generators} (\texttt{stair\_cert.py} and the search for the potential parameters) are untrusted; only acceptance of their output by the Lean checkers matters.
\item The scripts in \texttt{scripts/} use exact rational arithmetic, except \texttt{staircase.py} (floating point, and decimal interval arithmetic for the threshold) and the bisections for displayed roots; \texttt{replay39.py} and \texttt{w13\_25line.py} are written from Williams's text, \texttt{dag\_check.py} and \texttt{abs\_check.py} from the Lean definitions. None is part of the kernel-checked results.
\end{itemize}

\paragraph{(c) Scope.}
The results are about one proof calculus, the one of \cite[\S7]{W13} as formalised, with the variants of Section~\ref{sec:coverage}. They say nothing about lower bounds for SAT on these machines proved by other methods.

\section{Relation to earlier work}
For $d=1$, \cite[p.~A:31]{W13} recalls that Kannan proved $\NTIME[n]\not\subseteq\DTIME_1[n^{\sqrt[4]{3/2}}]$, and that the argument with Lemmas~7.1 and~3.3 gives $\sqrt{3/2}$, ``the SAT lower bound proved by Van Melkebeek and Raz'' \cite{vMR05}; for general $d$ the previous bound was $\sqrt{(d+2)/(d+1)}$ (p.~A:3). Williams \cite{W13} obtained $r_d$ (Theorem~7.5) and the limit $1+1/(d+1)$ for the calculus (Theorem~7.9). Theorem~\ref{thm:lower} improves $r_d$ to $1.30921174$, $1.1915968$, $1.1386523$, conditionally on the same lemmas, and Theorems~\ref{thm:ceiling} and~\ref{thm:sqrt} lower the limit of the formal calculus to $1.3094$, $1.192$, $1.1392$, and to $\sqrt{1+1/d}$ for every $d>0$. For the $\mathsf{DTS}$ setting of \cite[\S4]{W13}, Buss and Williams \cite{BW15} proved that $2\cos(\pi/7)$ is optimal for alternation-trading proofs.

\section*{Licence}
This text is licensed under CC BY 4.0. The Lean code and the scripts are licensed under the Apache License~2.0.

\begin{thebibliography}{Wfull}

\bibitem[BW15]{BW15}
S.~R. Buss and R.~Williams.
\newblock Limits on alternation trading proofs for time--space lower bounds.
\newblock \emph{Computational Complexity}, 24(3):533--600, 2015.
\newblock \href{https://doi.org/10.1007/s00037-015-0104-9}{doi:10.1007/s00037-015-0104-9}.

\bibitem[vMR05]{vMR05}
D.~van Melkebeek and R.~Raz.
\newblock A time lower bound for satisfiability.
\newblock \emph{Theoretical Computer Science}, 348(2--3):311--320, 2005.

\bibitem[Wfull]{Wfull}
R.~R. Williams.
\newblock Automated proofs of time lower bounds.
\newblock Long version of \cite{W13}, Carnegie Mellon University, \url{https://www.cs.cmu.edu/~ryanw/automated-lbs-full.pdf}. Its \S5 treats the hybrid model with $d=1$; we cite its Lemma~5.2 and the end of \S5.

\bibitem[W13]{W13}
R.~Williams.
\newblock Alternation-trading proofs, linear programming, and lower bounds.
\newblock \emph{ACM Transactions on Computation Theory}, 5(2):Article~6, 1--49, 2013.
\newblock \href{https://doi.org/10.1145/2493246.2493249}{doi:10.1145/2493246.2493249}.
\newblock Author copy: \url{https://people.csail.mit.edu/rrw/automated-lbs.pdf}, whose page numbers (A:1--A:40) we use.

\end{thebibliography}

\end{document}
